\section{Auxiliary validation and projection figures}
\label{app:validation-figures}

This appendix collects figures that are useful for review and Overleaf continuity but
are too detailed for the main PRD narrative. They should be pruned before submission
unless they become part of the quoted systematic uncertainty or a referee-facing
validation argument.

\begin{figure*}[t]
\centering
\ifigure{0.92\textwidth}{gpr_validation_flow.png}
\caption{Current HPS GPR bump-hunt and closure-validation flow. The diagram records
the analysis-note bookkeeping from spectrum preparation through GPR refits, signal
extraction, $\CLs$ limits, and functional-form closure tests.}
\label{fig:app-validation-flow}
\end{figure*}

\begin{figure*}[t]
\centering
\begin{minipage}{0.32\textwidth}
\centering
\ifigure{\linewidth}{closure_2015.png}\\
(a) 2015 closure suite.
\end{minipage}
\hfill
\begin{minipage}{0.32\textwidth}
\centering
\ifigure{\linewidth}{closure_2016.png}\\
(b) 2016 10\% closure suite.
\end{minipage}
\hfill
\begin{minipage}{0.32\textwidth}
\centering
\ifigure{\linewidth}{closure_2021.png}\\
(c) 2021 1\% closure suite.
\end{minipage}
\caption{Dataset-level full-refit closure summaries from the current validation
program. These panels diagnose pull means, pull widths, recovered amplitudes, and
significance residuals for the staged samples.}
\label{fig:app-closure-suites}
\end{figure*}

\begin{figure*}[t]
\centering
\ifigure{0.82\textwidth}{observed_extraction_119MeV.png}
\caption{Representative observed shared-coupling extraction display at a mass
hypothesis of 119 MeV. This figure is useful for checking how the data-set-level GPR
fits and signal templates enter the common $\eps^2$ likelihood, but it is not
necessary for the compact publication narrative.}
\label{fig:app-observed-display}
\end{figure*}

\begin{figure*}[t]
\centering
\ifigure{0.82\textwidth}{limit_tail_diagnostics.png}
\caption{Observed-limit tail diagnostics for the staged simultaneous analysis. These
curves compare the observed upper limit with background-only upper-limit ensembles.
They diagnose unusually strong or weak limits and are distinct from the discovery
$p_0$ statistic.}
\label{fig:app-tail-diagnostics}
\end{figure*}

\begin{figure*}[t]
\centering
\begin{minipage}{0.48\textwidth}
\centering
\ifigure{\linewidth}{cl90_vs_cl95.png}\\
(a) Combined 90\% and 95\% comparison.
\end{minipage}
\hfill
\begin{minipage}{0.48\textwidth}
\centering
\ifigure{\linewidth}{scaled_2021_resolution.png}\\
(b) 2021 scaled-resolution comparison.
\end{minipage}
\caption{Confidence-level and 2021 resolution systematic cross-checks. These studies
are important for freezing the final result convention but should remain auxiliary
unless they are folded into the quoted systematic model.}
\label{fig:app-systematic-crosschecks}
\end{figure*}

\begin{figure*}[t]
\centering
\begin{minipage}{0.48\textwidth}
\centering
\ifigure{\linewidth}{full_projection_vs_2016.png}\\
(a) Projected full-data reach versus the published 2016 HPS prompt limit.
\end{minipage}
\hfill
\begin{minipage}{0.48\textwidth}
\centering
\ifigure{\linewidth}{projection_vs_babar95.png}\\
(b) Projected full-data reach versus BaBar at an approximate 95\% confidence level.
\end{minipage}
\caption{Projection-only comparison figures. These plots are useful for planning the
eventual external-contour narrative, but they should not be described as observed
exclusions from the current staged samples.}
\label{fig:app-projections}
\end{figure*}
